\section{Wireless Communication and Zigbee Network}
\subsection{Wireless Communication Selection}
Wireless communication will be used in the application to transmit information from each module to the user. In each windmill app. 1000 sensor modules will be needed, one for each bolt and joint, therefore an appropriate solution would be to use a mesh network structure protocol as zigbee. A zigbee mesh network structure is build up of three components: The coordinator, the routers and the end devices. The end devices are each sensor module, the routers helps to make bridges and carry on the signal between end devices, and the coordinator collects data from all end devices and transmits it to the home assistant application. An advantage of the mesh network structure is that it is selfhealing, and can reach further than eg. a star formation, which only has communication between master and end device. A visual representation of a zigbee mesh network can be seen in \cref{fig:Zigbee_Network}.

\begin{figure}[H]
\centering
\includestandalone[mode=buildmissing, width=\linewidth]{Standalone/Drawings/wireless_com_selection-tikz}
    \caption{Zigbee mesh network structure and star network structure.}
    \label{fig:Zigbee_Network}
\end{figure}



Compared to similar network protocols such as Bluetooth Low Energy (BLE) and thread, Zigbee is preferable for this application, because it can host many devices and has relatively low complexity \cite{Thread_vs_zigbee, Zigbee_alliance}. BLE is a network protocol designed for low-power, short-range applications, but is more suited for peer-to-peer communcation or to be used in a star formation and Thread is a IPv6-based mesh networking protocol built on IEEE 802.15.4 at 2.4 GHz, where each nodes receives it's own IP adress, which enables direct IP communication and simplified integration with cloud services. Eventhough this seems preferable it requires additional layers such as 6LoWPAN, UDP, and IPv6 header handling, which makes a more complicated process with increased stack complexity, memory usage, and firmware size.

Zigbee uses a purpose-built, lightweight mesh networking layer with 16-bit short addressing managed by a coordinator. This results in smaller packet overhead and a simpler state machine for sleepy end devices. For sensors that wake up only once per month, this simplicity improves predictability and reduces implementation risk, especially in harsh offshore environments where maintenance access is limited.

From a hardware perspective, Zigbee and Thread can run on the same 802.15.4 SoCs. However, Zigbee stacks typically require slightly less flash and RAM, allowing the use of compact integrated solutions such as QFN packages around 5×5 mm (e.g., EFR32MG22, CC2652R). Using an integrated Zigbee SoC avoids the need for larger module-based Thread solutions or more powerful MCUs, which may increase PCB size and power consumption. 

In an offshore wind turbine environment, metal structures and turbine geometry can significantly attenuate and reflect RF signals, and no 2.4 GHz protocol inherently provides superior RF penetration under these conditions. Consequently, protocol selection should prioritize mesh reliability, implementation complexity, ecosystem maturity, and long-term operational stability. For applications where a mature and widely deployed mesh networking solution is preferred, Zigbee offers a proven ecosystem, established interoperability, and a long track record of industrial deployments, making it a strong candidate for this environment.

\subsection{Integrated vs. External Zigbee Module}
Each module requires a wireless communication circuit, which can either be integrated directly onto the module PCB or implemented using a separate, pre-certified RF module. The primary advantage of using an external RF module is the reduced design complexity, as the RF circuitry and antenna matching are already optimized and validated by the manufacturer. This can reduce development effort and lower the risk of RF performance issues. Additionally, it allows for a wider selection of microcontrollers, since wireless functionality does not need to be integrated into the MCU itself.

However, external RF modules generally increase both the cost and the physical size of the system. They may also introduce additional interfaces and integration considerations compared to a fully integrated solution.

After evaluating these trade-offs, the RF functionality was integrated directly into the system design. This approach reduces both the overall cost and PCB footprint while providing a more compact and efficient solution for the final product.

\subsection{The Zigbee Network}
\cref{fig:WC_Overview} illustrates the placements of the end notes, routers and the coordinator in a windmill. The yellow dots, representing the end devices, are to be placed on each nut around the flange, sending information to the nearest router. The routers are symbolized by the blue dots, which are spaced evenly throughout the windmill tower. The green dot represents the coordinator, which works as a host, collecting information from the end devices and transfers the data to a WiFi connection. For proof of concept each end device will include a 2.4 GHz antenna which will be used to transmit data to the routers and onward to the coordinator. For the coordinator a Sonoff Gateway ZigBee 3.0 USB Dongle Plus will be connected to a computer, which has been configured to work as a home assistant platform using the open-source Linux based system UBUNTU. 
The routers are often integrated into light bulbs or other analog outlets and will not be explored in the prototype setup.

\begin{figure}[H]
\centering
\includestandalone[mode=buildmissing, width=\linewidth]{Standalone/Drawings/WC_overview}
    \caption{Conceptual layout of wireless communication modules distributed throughout a wind turbine.}
    \label{fig:WC_Overview}
\end{figure}
